\section{March 19}


\subsection{Hamiltonian mechanics}

  Assume that $(M,g,J,\omega)$ is a (almost) K\"ahler manifold, where $g$ is the Riemannian metric, $J$ is the complex structure, and $\omega$ is the K\"ahler form.
  We have 
  \[ \omega(X,Y) = g(JX,Y), \quad \forall X,Y \in TM. \]
  For all $f \in C^\infty(M)$, let $\nabla f$ be the gradient of $f$ with respect to $g$.
  We have 
  \[ g \left( \nabla f, X \right) = X(f), \quad \forall X \in TM. \]
  Let $X_f$ be the Hamiltonian vector field of $f$ with respect to $\omega$, i.e.,
  \[ \omega(X_f, Y) = -Y(f), \quad \forall Y \in TM. \]
  Then we have
  \[ \omega(X_f, Y) = -g(\nabla f, Y) = g(J \nabla f, Y), \quad \forall Y \in TM. \] 
  Hence, we have $X_f = J \nabla f$.
  We prove the following lemma.

  \begin{lemma}
    Let $(M,g,J,\omega)$ be a (almost) K\"ahler manifold. 
    For all $f \in C^\infty(M)$, we have
    \[ X_f = J \nabla f. \]
  \end{lemma}

  \begin{example}
    Let $M = \bbR^{2n}$ with  real coordinates $(q_1, \ldots, q_n, p_1, \ldots, p_n)$.
    Then we have the standard K\"ahler structure $\omega = \sum_{i=1}^n dq_i \wedge dp_i$, $g(\partial_{q_i}, \partial_{q_j}) = g(\partial_{p_i}, \partial_{p_j}) = \delta_{ij}$, $g(\partial_{q_i}, \partial_{p_j}) = 0$ and $J \partial_{q_i} = \partial_{p_i}$, $J \partial_{p_i} = -\partial_{q_i}$.
    
    Let $H \in C^\infty(\bbR^{2n})$.
    Then we have
    \[ \nabla H = \sum_{i=1}^n \left( \frac{\partial H}{\partial q_i} \partial_{q_i} + \frac{\partial H}{\partial p_i} \partial_{p_i} \right). \]
    It follows that
    \[ X_H = J \nabla H = \sum_{i=1}^n \left( - \frac{\partial H}{\partial p_i} \partial_{q_i} + \frac{\partial H}{\partial q_i} \partial_{p_i} \right). \]
    Let $\gamma(t) = (q_1(t), \ldots, q_n(t), p_1(t), \ldots, p_n(t))$ be an integral curve of $X_H$, i.e., $\dot{\gamma}(t) = X_H(\gamma(t))$.
    Then we have
    \[ \dot{q}_i(t) = -\frac{\partial H}{\partial p_i}, \quad \dot{p}_i(t) = \frac{\partial H}{\partial q_i}, \quad i = 1, \ldots, n. \]
    These are the \emph{Hamilton's equations}.
    % \Yang{}
  \end{example}

  \begin{example}
    Assume a particle of mass $m$ moves in configuration space $\bbR^3$ under a potential $V \in C^\infty(\bbR^3)$.
    If the trajectory of the particle is $\gamma(t) = (q_1(t), q_2(t), q_3(t))$, by Newton's second law, we have
    \[ m \ddot{\gamma}(t) = -\nabla V. \]
    In coordinates, we have
    \[ m \ddot{q}_i(t) = -\frac{\partial V}{\partial q_i}, \quad i = 1,2,3. \]
    Set $p_i = m \dot{q}_i$.
    This is the \emph{momentum} of the particle.
    Then we have
    \begin{equation}\label{eq:Newton_second_law}
      \dot{p}_i(t) = m \ddot{q}_i(t) = -\frac{\partial V}{\partial q_i}, \quad i = 1,2,3.
    \end{equation}  
    Set $\bbR^6 = T^*\bbR^3$ with coordinates $(q_1, q_2, q_3, p_1, p_2, p_3)$.
    This is the \emph{phase space} of the particle.

    Let 
    \[ H = -\left( \frac{1}{2m} \|p\|^2 + V(q) \right) \in C^\infty(\bbR^6). \]
    The $-H$ is the \emph{total energy} of the particle.
    Let $\beta(t) = (q_1(t), q_2(t), q_3(t), p_1(t), p_2(t), p_3(t))$ be an integral curve of $X_H$.
    We have 
    \[ \frac{\partial H}{\partial q_i} = -\frac{\partial V}{\partial q_i}, \quad \frac{\partial H}{\partial p_i} = -\frac{1}{m} p_i, \quad i = 1,2,3. \]
    By Hamilton's equations, we have
    \begin{equation}\label{eq:hamilton_equations}
      \dot{q}_i(t) = -\frac{\partial H}{\partial p_i} = \frac{1}{m} p_i, \quad \dot{p}_i(t) = \frac{\partial H}{\partial q_i} = -\frac{\partial V}{\partial q_i}, \quad i = 1,2,3.
    \end{equation}
    Hence the Hamilton's equations \eqref{eq:hamilton_equations} are equivalent to Newton's second law \eqref{eq:Newton_second_law}.
  \end{example}

  Let $(M,\omega)$ be a symplectic manifold, $H \in C^\infty(M)$ and $X_H$ be the Hamiltonian vector field of $H$.
  Let $\gamma(t)$ be an integral curve of $X_H$.
  
  \begin{lemma}
    A function $f \in C^\infty(M)$ is a constant of motion, i.e., $f(\gamma(t))$ is constant for all $t$, if and only if $\{H,f\} = 0$.
  \end{lemma}
  \begin{proof}
    Let $\gamma(t)$ be an arbitrary integral curve of $X_H$.
    We have
    \[ \frac{d}{dt} f(\gamma(t)) = df(\dot{\gamma}(t)) = df(X_H) = -i_{X_f} \omega(X_H) = -\omega(X_f, X_H) = -\{f,H\}|_{\gamma(t)}. \]
    The lemma follows immediately.
  \end{proof}

  \begin{definition}
    If $\{f,H\} = 0$, we say that $f$ is a \emph{first integral} or \emph{integral of motion} of $H$.
  \end{definition}

  Note that $H$ is always a first integral of itself since $\{H,H\} = 0$.
  This is the \emph{conservation of energy}.


\subsection{Darboux's theorem}

  \begin{theorem}\label{thm:Darboux}
    For any symplectic manifold $(M,\omega)$ and any point $p \in M$, there exist local coordinates $(x_1, \ldots, x_m, y_1, \ldots, y_m)$ around $p$ such that
    \[ \omega = \sum_{i=1}^m dx_i \wedge dy_i. \]
   In particular, any two symplectic manifolds of the same dimension are locally symplectomorphic.
  \end{theorem}
  \begin{proof}
    Induct on $m$.
    Assume that the theorem holds for $m-1$.
    Take $y_1 \in C^\infty(M)$ such that $dy_1(p) \neq 0$.
    Let $X$ be the Hamiltonian vector field of $-y_1$.
    Then we have $i_X \omega = dy_1$.
    Then $X(p) \neq 0$ since $dy_1(p) \neq 0$.
    There exists a function $x_1$ defined in a neighborhood of $p$ such that $X(x_1) \equiv 1$ on this neighborhood.
    Without loss of generality, replace $M$ by this neighborhood and assume that $X(x_1) \equiv 1$ on $M$.
    Let $Y$ be the Hamiltonian vector field of $x_1$.
    Then we have $i_Y \omega = -dx_1$.
    We have
    \[ \omega(X,Y) = -(i_Y \omega)(X) = dx_1(X) \equiv 1. \]
    By \cref{...}, since $X,Y$ are symplectic vector fields, we have 
    \[ i_{[X,Y]} \omega =  - d(\omega(X,Y)) = 0. \]
    Hence, we have $[X,Y] = 0$ by the non-degeneracy of $\omega$.
    Then 
    \[ Y(y_1) = dy_1(Y) = (i_X \omega)(Y) = \omega(X,Y) \equiv 1. \]
    By Frobenius' theorem, there exist local coordinates $(x_1, y_1, z_1, \ldots, z_{2m-2})$ around $p$ such that $X = \partial_{x_1}$ and $Y = \partial_{y_1}$.
    Let $\omega' = \omega - dx_1 \wedge dy_1$.
    Then $\omega'$ is a closed 2-form.
    We have $i_X \omega' = i_X \omega - i_X (dx_1 \wedge dy_1) = dy_1 - dy_1 = 0$ and $i_Y \omega' = i_Y \omega - i_Y (dx_1 \wedge dy_1) = -dx_1 + dx_1 = 0$.
    So $\omega' = \sum_{i=1}^{2m-2} a_{ij} dz_i \wedge dz_j$.
    Since $d\omega' = 0$, we have $a_{ij}$ does not depend on $x_1$ and $y_1$.
    So $\omega'$ can be viewed as a closed 2-form on the $(2m-2)$-dimensional manifold with coordinates $(z_1, \ldots, z_{2m-2})$.
    Note that 
    \[ \wedge^m \omega = \left(\wedge^{m-1} \omega'\right) \wedge dx_1 \wedge dy_1 \neq 0. \]
    Hence, $\wedge^{m-1} \omega' \neq 0$.
    It follows that $\omega'$ is a symplectic form on the $(2m-2)$-dimensional manifold with coordinates $(z_1, \ldots, z_{2m-2})$.
    By the induction hypothesis, there exist local coordinates $(x_2, \ldots, x_m, y_2, \ldots, y_m)$ around $p$ such that $\omega' = \sum_{i=2}^m dx_i \wedge dy_i$.
    Hence, we have $\omega = \sum_{i=1}^m dx_i \wedge dy_i$.
  \end{proof}

  \begin{corollary}
    There is no local invariant of symplectic manifolds other than the dimension.
  \end{corollary}

  Given any symplectic manifold $(M,\omega)$ and $H \in C^\infty(M)$.
  Let $f$ be a first integral of $H$, i.e., $\{H,f\} = 0$.
  For $p \in M$ with $df(p) \neq 0$, by the proof of Darboux's theorem (\cref{thm:Darboux}), there exist local coordinates $(x_1, \ldots, x_m, y_1, \ldots, y_m)$ around $p$ such that $\omega = \sum_{i=1}^m dx_i \wedge dy_i$ and $f = y_m$.
  With the standard K\"ahler structure with respect to these coordinates, we have 
  \[ X_f = J \nabla f = -\partial_{x_m}. \]
  Since $\{H,f\} = 0$, we have 
  \[ \frac{\partial H}{\partial x_m} = -\{H,f\} = 0. \]
  Hence, $H$ does not depend on $x_m$.

